\subsection{Сохранение нагрузочных зависимостей после discovery}
Все 288 обучений завершены. Каждый калибровочный и валидационный эпизод при нагрузочных возвратах удовлетворяет предкризисному критерию. Один из 3168 эпизодов, в P4 при combined feedback, нарушает его; этот набор сохранён в диагностических результатах.

При автоматически выбранной структуре retained-load снижает среднюю нормированную стоимость с равными весами процессов с 0.627 до 0.542, на 13.5\% относительно предиктивного отбора. Основная парная разность равна $-0.0848$, скорректированный 97.5\%-ный интервал — $[-0.1353,-0.0343]$. Таблица~\ref{tab:structures_results} показывает отдельные структуры и сильные альтернативы. В исходных шкалах средняя стоимость снижается в P1, P2, P3, P5 и P6 соответственно на 11.6\%, 14.5\%, 7.0\%, 31.0\% и 24.9\%. После поправки на шесть вторичных сравнений преимущество P6 сохраняется: интервал разности составляет $[-18.48,-1.82]$ тыс. В P4 стоимость, напротив, увеличивается на 16.8\%, скорректированный интервал — $[2.30,8.74]$ тыс. 

\begin{table}[t]\caption{Routing feedback при выбранной обнаруженной структуре: стоимость каждого обученного варианта нормирована на соответствующую стоимость без вмешательства. Regret равен разности стоимости retained-load и планировщика с известным генератором, также нормированной на no action. Каждая строка содержит 12 независимых обучений.}\label{tab:structures_results}\centering\small\setlength{\tabcolsep}{4pt}
\begin{tabular}{lrrrrrr}\toprule
Процесс & Static & Predictive & Retained & Heuristic & Generator & Regret\\\midrule
P1 & 0.584 & 0.555 & 0.491 & 0.499 & 0.429 & 0.062 \\
P2 & 0.688 & 0.658 & 0.561 & 0.648 & 0.435 & 0.126 \\
P3 & 0.659 & 0.659 & 0.613 & 0.635 & 0.522 & 0.092 \\
P4 & 0.518 & 0.527 & 0.615 & 0.623 & 0.513 & 0.102 \\
P5 & 1.030 & 0.828 & 0.571 & 0.467 & 0.415 & 0.155 \\
P6 & 0.573 & 0.537 & 0.404 & 0.574 & 0.351 & 0.054 \\
\bottomrule\end{tabular}\end{table}


Контрольные механизмы уточняют условия пользы от сохранения компонентов. Описательная разность нормированной стоимости retained минус predictive равна $-0.055$ при combined feedback, но $+0.050$ при фиксированном поведении и $+0.110$ при service feedback. Сохранение слабо оценённых зависимостей возвратов полезно в routing-сравнении, но может вносить ошибку при отсутствии этих зависимостей. Рисунок~\ref{fig:structural_matrix} разделяет механизмы и две процедуры discovery.

\begin{figure}[t]\centering
\includegraphics[width=\textwidth]{generated/structural_matrix_ru.pdf}
\caption{Среднее снижение стоимости retained-load относительно предиктивного отбора внутри каждого структурного варианта, в процентах. Положительные значения соответствуют преимуществу сохранения зависимостей. Обе процедуры discovery используют один исправленный исполнитель. Оценки ячеек описательные; основные агрегированные сравнения приведены в тексте.}\label{fig:structural_matrix}
\end{figure}

\subsection{Исполнимая структура сохраняет смысл вмешательства}
Во всех 288 обучениях выбранные обнаруженные модели имеют тот же язык видимых последовательностей, что и исходный процесс; средняя обучающая precision составляет от 0.996 до 1.000. Они точно воспроизводят планировочные стоимости при известной структуре для всех четырёх комбинаций компонентов длительности и маршрутизации. Стандартный IM восстанавливает исходный язык в 144 случаях; P1, P2 и P6 допускают дополнительные пути. Во втором основном сравнении переход от стандартного IM к выбранной структуре снижает нормированную стоимость retained-load с 0.638 до 0.542, на 15.0\%. Парная разность равна $-0.0958$, скорректированный 97.5\%-ный интервал — $[-0.1205,-0.0711]$.

P2 показывает последствие для выбора реакции. Выбранная структура приводит к сокращению rework в 8 из 12 обучений, а стандартный IM — к увеличению мощности во всех случаях; нормированные стоимости составляют 0.561 и 0.898. Оба варианта используют одинаковое обученное поведение и один исправленный исполнитель. Сохранение структуры возвратов и наблюдаемых условий маршрутизации меняет действие, которое поддерживают эти оценки поведения. Одного воспроизведения обучающих трасс недостаточно для проверки пригодности модели к планированию вмешательств.

\subsection{Системное планирование и простые правила}
P6 при нагрузочных возвратах даёт диагностический пример преимущества симуляции перед локальной эвристикой: нормированные стоимости равны 0.404 и 0.574, описательный парный 95\%-ный интервал разности — $[-24.85,-0.47]$ тыс. Эвристика увеличивает мощность отдельной ветви в 4 обучениях; retained-load выбирает сокращение возвратов или общий ресурс. Сравнения P3 и P5 не устанавливают преимущества над эвристикой. При service feedback в P6 retained-load также имеет меньшую среднюю стоимость, 0.671 против 0.878, но статическое планирование стоит лишь 0.496. Симуляция способна помочь оценить системные последствия без необходимости сохранять каждый нагрузочный компонент.

Проверка ресурсов показывает, почему эти последствия существенны. В P5 при combined feedback дополнительный работник B снижает среднюю очередь с 137.95 до 37.29, одновременно увеличивая очередь D с 1.28 до 116.13; доминирующая очередь перемещается во всех 24 парных диагностических траекториях. Однако точная структура не гарантирует точной экстраполяции поведения. В P4 при нагрузочных возвратах удаление обученных нагрузочных членов маршрутизации снижает нормированную стоимость с 0.615 до 0.514; удаление членов длительности оставляет её на уровне 0.615. Эта абляция локализует ошибку решения в оценённых зависимостях возвратов.

\subsection{Чувствительность к отбору и эффектам вмешательства}
В P1, P2, P3, P5 и P6 retained-load имеет меньшую среднюю routing-стоимость при каждой проверенной настройке любого из двух порогов селектора. Изменение $\Delta_R$ от 0 до 0.01 меняет нормированную стоимость предиктивного отбора в P5 с 0.726 до 0.928, тогда как retained-load стоит 0.571. В P4 предиктивный отбор дешевле на всей сетке. Эффект выходит за пределы исходного порога селектора, оставаясь зависимым от процесса. Перебор диагностирует фиксированные методы, а не выбирает новый порог по тесту.

При увеличении доли сложных экземпляров P2 до 65\% неизменённые политики retained и predictive стоят 0.668 и 0.679. В P3 фактическое сокращение rework на 30\% даёт стоимости 0.730 и 0.683, а на 70\% — 0.518 и 0.639. Порядок меняется без повторного выбора политики. Поэтому подготовленная реакция должна сопровождаться явным диапазоном эффектов вмешательства, в котором её рекомендация остаётся полезной.
